% ============================================================
% IJHSR REVIEW ARTICLE MOCKUP
% International Journal of High School Research
% Engineering and Computer Science
%
% Compile with XeLaTeX or LuaLaTeX
% ============================================================

\documentclass[10pt]{article}

% ============================================================
% COMPILER CHECK
% ============================================================

\usepackage{iftex}

\ifPDFTeX
  \errmessage{This document requires XeLaTeX or LuaLaTeX}
\fi

% ============================================================
% PAGE SETUP
% ============================================================

\usepackage[
    letterpaper,
    top=1in,
    bottom=1in,
    left=1in,
    right=1in
]{geometry}

% ============================================================
% FONT
% IJHSR specifies Arial.
% If Arial is unavailable in the LaTeX environment,
% TeX Gyre Heros is used as a visual fallback.
% ============================================================

\usepackage{fontspec}

\IfFontExistsTF{Arial}
  {\setmainfont{Arial}}
  {\setmainfont{TeX Gyre Heros}}

% ============================================================
% BODY FORMATTING
% ============================================================

\usepackage{setspace}
\singlespacing

\usepackage{ragged2e}
\justifying

\setlength{\parindent}{0pt}
\setlength{\parskip}{6pt}
\setlength{\emergencystretch}{2em}

\usepackage{microtype}

% ============================================================
% SECTION FORMATTING
% ============================================================

\usepackage{titlesec}

\setcounter{secnumdepth}{0}

% Major headings:
% Arial 14 pt, bold, 8 pt after
\titleformat{\section}
  {\fontsize{14}{14}\selectfont\bfseries}
  {}
  {0pt}
  {}

\titlespacing*{\section}
  {0pt}
  {10pt}
  {8pt}

% Subheadings:
% Arial 12 pt, bold italic, 8 pt after
\titleformat{\subsection}
  {\fontsize{12}{12}\selectfont\bfseries\itshape}
  {}
  {0pt}
  {}

\titlespacing*{\subsection}
  {0pt}
  {10pt}
  {8pt}

% ============================================================
% MATH
% ============================================================

\usepackage{amsmath}
\usepackage{amssymb}

% ============================================================
% TABLES
% ============================================================

\usepackage{booktabs}
\usepackage{tabularx}
\usepackage{array}
\usepackage{float}

\renewcommand{\arraystretch}{1.18}
\setlength{\tabcolsep}{4pt}

\newcolumntype{Y}{>{\RaggedRight\arraybackslash}X}
\newcolumntype{L}[1]{>{\RaggedRight\arraybackslash}p{#1}}

% ============================================================
% FIGURES AND CAPTIONS
% ============================================================

\usepackage{graphicx}
\usepackage{caption}

\DeclareCaptionFont{ijhsr}{
    \fontsize{11}{12}\selectfont
}

\captionsetup{
    font=ijhsr,
    justification=RaggedRight,
    singlelinecheck=false,
    skip=8pt
}

% Table captions above
\captionsetup[table]{
    position=above
}

% Figure captions below
\captionsetup[figure]{
    position=below
}

% ============================================================
% LINKS
% ============================================================

\usepackage[hidelinks]{hyperref}
\urlstyle{same}

% ============================================================
% LISTS
% ============================================================

\usepackage{enumitem}

\setlist[itemize]{
    leftmargin=1.5em,
    itemsep=2pt,
    topsep=3pt
}

\setlist[enumerate]{
    leftmargin=1.7em,
    itemsep=2pt,
    topsep=3pt
}

% ============================================================
% SUPERSCRIPT ACS-STYLE CITATION NUMBERS
% Citation numbers should appear AFTER punctuation.
%
% Example:
% Numerical reliability is important.\supercite{1}
% ============================================================

\newcommand{\supercite}[1]{\textsuperscript{#1}}

% ============================================================
% DOCUMENT
% ============================================================

\begin{document}

% ============================================================
% TITLE
% 16 pt, bold, title case, maximum two lines
% ============================================================

{\fontsize{16}{17}\selectfont\bfseries
Reliability of Transonic Airfoil CFD: A Review of Mesh Resolution,
Turbulence Modeling, and Shock Prediction
\par}

\vspace{8pt}

% ============================================================
% AUTHOR INFORMATION
% ============================================================

Adhvaith Ananth (student)

Texas Academy of Mathematics and Science,
University of North Texas,
1155 Union Circle \#305309,
Denton, Texas 76203-5017, USA

% Personal contact email omitted from the public portfolio copy.

\vspace{8pt}

% ============================================================
% ABSTRACT
% Maximum 200 words
% One paragraph
% No citations
% ============================================================

\section{Abstract}

Computational fluid dynamics (CFD) is widely used to study aerodynamic
flows, but a numerical solution is useful only when it represents the
underlying physics with sufficient accuracy. This challenge is
particularly important in transonic aerodynamics, where localized
supersonic regions, shock waves, turbulent boundary layers, and possible
flow separation interact. This review examines major factors that
influence the reliability of transonic airfoil CFD, including mesh
resolution, near-wall treatment, turbulence modeling, convergence
assessment, and experimental validation. The RAE 2822 airfoil is used
as a principal benchmark because extensive experimental pressure and
boundary-layer measurements are available for numerical comparison.
Published studies demonstrate that differences in operating conditions,
grid design, turbulence closure, and numerical methods can alter shock
prediction and aerodynamic coefficients. Shock-control bumps are
considered as an application in which errors in predicted shock position
can directly influence aerodynamic design. Overall, reliable transonic
CFD requires multiple forms of evidence rather than reliance on cell
count, residual reduction, turbulence-model choice, or visually
plausible contours alone.

% ============================================================
% KEYWORDS
% IJHSR: maximum 5
% Keyword 1 = broad field
% Keyword 2 = subfield
% Comma-separated
% ============================================================

\section{Keywords}

Engineering,
Aerospace and Aeronautical Engineering,
Computational Fluid Dynamics,
Transonic Aerodynamics,
Shock Control Bumps

% ============================================================
% INTRODUCTION
% ============================================================

\section{Introduction}

Computational fluid dynamics has become one of the most widely used
numerical tools in aerospace engineering. CFD allows engineers to
approximate the governing equations of fluid motion throughout a
computational domain and examine quantities including pressure,
velocity, density, temperature, aerodynamic loads, and flow structures.

The ability to calculate a numerical flow field, however, does not
automatically establish that the result accurately represents the
physical flow. A CFD solution can appear plausible while still being
affected by discretization error, insufficient convergence, incorrect
boundary conditions, inadequate mesh resolution, or limitations in the
physical models used by the solver.

Oberkampf and Trucano distinguish between verification and validation
in computational modeling.\supercite{1} Verification asks whether the
mathematical model is being solved correctly, while validation asks
whether that model represents physical reality with sufficient accuracy
for its intended application. A calculation can therefore be
numerically converged without necessarily being physically accurate.

Transonic aerodynamics presents a demanding environment for numerical
simulation. Although the freestream Mach number can remain below one,
flow over an airfoil may accelerate locally beyond the speed of sound.
The return to subsonic conditions commonly occurs through a shock wave,
producing a rapid pressure increase that can interact with the turbulent
boundary layer.

Depending on the operating condition, shock-wave/boundary-layer
interaction can increase drag, alter lift, thicken or separate the
boundary layer, and contribute to unsteady aerodynamic buffet.
Predicting this behavior requires a numerical method that can represent
both compressible shock physics and viscous near-wall flow.

The RAE 2822 airfoil has become an important benchmark for transonic CFD
because experimental surface-pressure, boundary-layer, and wake
measurements are available for comparison.\supercite{2} This review
uses the RAE 2822 and related studies to examine how mesh construction,
turbulence modeling, numerical convergence, and validation methodology
influence the reliability of transonic CFD.

Shock-control bumps are also considered as an application of this
problem. These localized surface modifications are designed to alter
the compression process near a transonic shock. Their effectiveness
depends strongly on the predicted shock location, making them a useful
example of why numerical validation should precede CFD-based
optimization.

% ============================================================
% RESULTS AND DISCUSSION
% Required IJHSR major heading
% ============================================================

\section{Results and Discussion}

% ============================================================
% LITERATURE SEARCH
% ============================================================

\subsection{Literature Selection and Scope}

This article is a targeted narrative review rather than a systematic
review. Literature was identified through focused searches of the NASA
Glenn Research Center NPARC Alliance Validation Archive and scholarly
publisher and DOI databases. Search terms included combinations of
``RAE 2822 transonic validation,'' ``CFD verification validation,''
``shock wave boundary layer turbulence model,'' ``transonic airfoil
mesh convergence,'' and ``shock control bump review.''

Sources were selected when they addressed at least one of four areas:
experimental or numerical RAE 2822 benchmarking, CFD verification and
discretization uncertainty, turbulence modeling for
shock-wave/boundary-layer interaction, or shock-control-bump
performance. Reference chaining from major review and application
studies was also used to identify foundational publications.

Preference was given to studies that described operating conditions,
numerical methods, validation evidence, or experimental comparisons in
sufficient detail to permit meaningful interpretation.

Because the search was targeted rather than systematic, the literature
discussed here should not be interpreted as an exhaustive collection of
all published work on transonic CFD.

% ============================================================
% TRANSONIC PHYSICS
% ============================================================

\subsection{Physics of Transonic Airfoil Flow}

The local Mach number is defined as

\begin{equation}
M=\frac{V}{a},
\end{equation}

where $V$ is local flow velocity and $a$ is the local speed of sound.

A freestream Mach number below one does not require the complete flow
around an airfoil to remain subsonic. Pressure gradients created by the
airfoil can accelerate the flow, particularly over the upper surface.
Under transonic conditions this can create a localized region in which
the Mach number exceeds unity.

The downstream transition from local supersonic flow to subsonic flow
can occur through a shock wave. Across the shock, pressure, density, and
temperature increase while velocity decreases.

In a CFD calculation the discontinuity is represented across a finite
number of computational cells. An insufficiently resolved shock can
therefore become numerically smeared or appear at an incorrect
streamwise location.

The pressure increase through the shock also produces a strong adverse
pressure gradient. Fluid close to the airfoil surface has less momentum
than the outer flow, allowing the boundary layer to thicken or separate
when exposed to a sufficiently strong pressure rise.

This interaction is known as shock-wave/boundary-layer interaction
(SWBLI). Turbulence modeling and separation prediction can strongly
influence its numerical representation.\supercite{3}

% ============================================================
% VERIFICATION AND VALIDATION
% ============================================================

\subsection{Verification, Validation, and Numerical Uncertainty}

Verification and validation address different aspects of computational
reliability.

Verification examines whether the numerical implementation solves the
intended equations correctly. Validation examines whether the physical
model reproduces experimental reality with sufficient accuracy for the
intended application.\supercite{1}

This distinction is important because the governing equations are
approximated on a finite grid. Discretization introduces numerical
error, particularly in regions containing large gradients such as
boundary layers and shock waves.

One common method of assessing discretization sensitivity is to compare
solutions generated on progressively refined grids. Quantities such as
pressure coefficient, shock location, lift, drag, wall shear, and
separation position can then be compared.

A simple difference measure between coarse and fine solutions can be
written as

\begin{equation}
\Delta Q =
\left|
\frac{Q_{\mathrm{fine}}-Q_{\mathrm{coarse}}}
{Q_{\mathrm{fine}}}
\right|\times100\%.
\end{equation}

More formal methods, including the Grid Convergence Index, provide
structured approaches for estimating uncertainty associated with
spatial discretization.\supercite{4}

Grid convergence should be evaluated for the quantities that matter to
the study. Lift may become relatively insensitive to grid refinement
while drag or shock location remains sensitive.

% ============================================================
% RAE 2822
% ============================================================

\subsection{RAE 2822 as a Validation Benchmark}

The RAE 2822 airfoil is widely used in transonic CFD validation because
of the experimental database reported by Cook, McDonald, and
Firmin.\supercite{2}

The dataset includes surface-pressure, boundary-layer, and wake
measurements. These measurements allow computational results to be
evaluated against physical data rather than solely against another
numerical prediction.

NASA's NPARC Alliance Validation Archive contains RAE 2822 cases
designed specifically to compare computational predictions with
experimental measurements.\supercite{5}

One NASA validation study uses a structured C-grid with dimensions of
369 by 65 points at a freestream Mach number of 0.729, an angle of
attack of approximately $2.31^\circ$, and a Reynolds number of roughly
$6.5\times10^6$.\supercite{5}

Multiple turbulence models are evaluated, and computed surface-pressure
coefficients are compared with experimental data.

Malkocoglu and coauthors used a different RAE 2822 case as a validation
step before analyzing shock-wave/boundary-layer interaction and
shock-control bumps on the NACA 64-A-307 airfoil.\supercite{6}

Their RAE configuration used Mach 0.75, an angle of attack of
$3.19^\circ$, and Reynolds number $6.2\times10^6$ with a modified
O-type mesh containing 545,768 cells.

These studies should not be treated as identical simply because both
use the RAE 2822 geometry. Their operating conditions, meshes, numerical
solvers, and turbulence-model comparisons differ.

Malkocoglu and coauthors reported better shock-location capture with
k--omega SST than with Spalart--Allmaras for their numerical
configuration.\supercite{6}

This is evidence of model sensitivity within that particular
configuration rather than proof that one turbulence model is universally
superior.

% ============================================================
% MESH
% ============================================================

\subsection{Mesh Resolution and Grid Topology}

A CFD mesh determines the spatial locations at which the governing
equations are approximated.

Transonic-airfoil simulations contain several regions where additional
resolution may be important: the leading edge, near-wall boundary layer,
expected shock region, trailing edge, and downstream wake.

C-type and O-type structured grids are common in benchmark airfoil CFD
because they allow direct control of surface-normal spacing and
streamwise cell distribution.

NASA's RAE 2822 validation uses a C-grid, whereas the Malkocoglu
validation employs a modified O-type grid.\supercite{5--6}

The existence of successful studies using different topologies shows
that grid topology alone does not determine accuracy. Local resolution,
cell quality, numerical formulation, and placement of refinement are
also important.

Unstructured meshes provide greater geometric flexibility and are
common for complex configurations. They should not automatically be
considered less accurate than structured grids. Their performance also
depends on local refinement, element quality, near-wall construction,
and numerical scheme.

Near the wall, resolution is commonly described using

\begin{equation}
y^+ = \frac{u_\tau y}{\nu},
\end{equation}

where $u_\tau$ is friction velocity, $y$ is wall-normal distance, and
$\nu$ is kinematic viscosity.

Wall-resolved approaches often target $y^+$ values close to unity.
However, a favorable first-cell distance alone cannot demonstrate mesh
adequacy. Streamwise spacing through the shock and resolution of the
wake remain important.

Raw cell count is therefore a weak basis for comparing unrelated CFD
studies. A larger grid does not automatically produce a more accurate
solution if its resolution is distributed inefficiently or if other
modeling errors dominate.

% ============================================================
% TURBULENCE
% ============================================================

\subsection{Turbulence Modeling}

Direct numerical simulation resolves turbulent motion across the
relevant range of scales but is computationally impractical for many
high-Reynolds-number engineering problems.

Large-eddy simulation reduces this requirement by modeling smaller
scales but remains computationally demanding for wall-bounded
aerodynamic flows.

Reynolds-averaged Navier--Stokes methods instead solve for averaged flow
quantities and model the influence of turbulent fluctuations. RANS
therefore remains common in engineering CFD.

The Spalart--Allmaras model is a one-equation turbulence model developed
for aerodynamic flows.\supercite{7}

Its comparatively low computational cost has contributed to its
widespread use in external aerodynamics.

Menter's k--omega shear-stress-transport model is a two-equation
eddy-viscosity formulation.\supercite{8}

Its treatment of near-wall flow and adverse pressure gradients has made
it common in simulations involving possible separation.

Comparative shock/boundary-layer studies show that turbulence closure
can materially alter predicted shock position and separation
behavior.\supercite{3}

Consequently, disagreement between CFD and experiment cannot
automatically be attributed to mesh resolution alone.

% ============================================================
% SHOCK EVALUATION
% ============================================================

\subsection{Evaluating Shock Prediction and Convergence}

Mach-number contours are useful for identifying supersonic regions and
visualizing shock structure, but they remain primarily qualitative.

Contour spacing, interpolation, and plotting choices can make a
numerical shock appear convincing even when its predicted position is
incorrect.

The surface pressure coefficient is

\begin{equation}
C_p =
\frac{p-p_{\infty}}
{\frac{1}{2}\rho_{\infty}V_{\infty}^{2}}.
\end{equation}

Pressure coefficient is especially useful for validation because it can
be measured experimentally and calculated numerically.

Shock location is often represented using normalized chord position
$x/c$. Because shock location responds to geometry, Mach number,
boundary-layer state, and turbulence modeling, it provides a sensitive
metric for comparing CFD with experiment.

Integrated aerodynamic forces provide additional information.

Lift and drag coefficients are

\begin{equation}
C_L=
\frac{L}
{\frac{1}{2}\rho_{\infty}V_{\infty}^{2}S},
\end{equation}

and

\begin{equation}
C_D=
\frac{D}
{\frac{1}{2}\rho_{\infty}V_{\infty}^{2}S}.
\end{equation}

Integrated coefficients, however, can hide local errors because
differences in surface pressure may partially cancel.

Agreement in $C_L$ alone therefore does not demonstrate that shock
position or surface loading has been accurately predicted.

Residual histories should also be considered alongside engineering
quantities.

Persistent variation in a CFD result may represent incomplete
convergence, numerical instability, or genuine physical unsteadiness.

% ============================================================
% SHOCK CONTROL BUMPS
% ============================================================

\subsection{Shock-Control Bumps as a Validation-Dependent Application}

Shock-control bumps are localized surface modifications intended to
alter the compression process associated with a transonic shock.

Bruce and Colliss reviewed shock-control-bump research and emphasized
the importance of bump shape, position, operating condition, and shock
behavior.\supercite{9}

A bump can alter the local surface curvature and compression-wave
structure, potentially weakening or redistributing the main shock.

Possible benefits include reduced wave drag and changes in separation
or buffet behavior.

The effectiveness of the geometry is strongly dependent on its position
relative to the shock. This creates a direct dependence on the accuracy
of the baseline CFD prediction.

If the baseline numerical method predicts the shock too far upstream or
downstream, an optimization algorithm may design a bump around a
numerical error rather than the physical shock position.

Tian and coauthors analyzed two shock-control-bump concepts on the
RAE 2822 using RANS and URANS calculations.\supercite{10}

Their results illustrate that a configuration that performs favorably
near one operating condition can provide less benefit at another.

The study therefore illustrates two related challenges: CFD must first
predict the baseline shock accurately, and the resulting geometry must
then be tested away from the single condition for which it was
optimized.

% ============================================================
% TABLE 1
% ============================================================

Table 1 compares several studies discussed throughout this review.
Differences in geometry, operating point, solver, and validation
evidence limit direct numerical comparison, but the studies collectively
show how different components of CFD reliability are evaluated.

\begin{table}[H]

\caption{\textbf{Comparison of selected transonic CFD and shock-control
studies.} The studies use different geometries, operating conditions,
and numerical methods; their numerical results therefore should not be
compared as identical test cases.}

\centering
\fontsize{9}{10}\selectfont

\begin{tabularx}{\textwidth}{
    L{2.0cm}
    L{2.15cm}
    L{2.15cm}
    L{2.4cm}
    Y
}

\toprule

\textbf{Study} &
\textbf{Configuration} &
\textbf{Method} &
\textbf{Validation evidence} &
\textbf{Key limitation or relevance} \\

\midrule

NASA RAE Study &
RAE 2822; $M=0.729$; $\alpha\approx2.31^\circ$ &
369 $\times$ 65 C-grid; multiple RANS models &
Experimental surface $C_p$ and aerodynamic convergence monitoring &
Strong benchmark evidence, but represents a specific RAE operating
point and numerical setup. \\

\addlinespace

Malkocoglu et al. &
RAE 2822; $M=0.75$; $\alpha=3.19^\circ$;
$Re=6.2\times10^6$ &
Modified O-grid; 545,768 cells; turbulence-model comparison &
Published RAE benchmark pressure comparison &
Useful validate-first workflow before NACA bump analysis; not directly
identical to the NASA RAE operating point. \\

\addlinespace

Liou et al. &
Multiple transonic and supersonic SWBLI cases &
Several turbulence models &
Comparison with experimental interaction data &
Demonstrates turbulence-model sensitivity but uses interaction cases
rather than one airfoil benchmark. \\

\addlinespace

Tian et al. &
RAE 2822 with two shock-control concepts &
RANS and URANS &
Baseline comparison and aerodynamic/buffet analysis &
Shows design and off-design sensitivity; conclusions depend on baseline
numerical fidelity. \\

\bottomrule

\end{tabularx}

\end{table}

% ============================================================
% TABLE 2
% ============================================================

Table 2 summarizes major sources of uncertainty that can influence
transonic-airfoil CFD.

\begin{table}[H]

\caption{\textbf{Major sources of uncertainty in transonic-airfoil
CFD.} No single assessment establishes numerical reliability; several
forms of evidence should be considered together.}

\centering
\fontsize{9}{10}\selectfont

\begin{tabularx}{\textwidth}{
    L{2.6cm}
    Y
    Y
}

\toprule

\textbf{Source} &
\textbf{Possible effect} &
\textbf{Assessment approach} \\

\midrule

Mesh density &
Shock smearing, pressure-gradient error, and altered force coefficients &
Grid-refinement comparison of $C_p$, $x/c$, $C_L$, and $C_D$. \\

\addlinespace

Near-wall spacing &
Incorrect boundary-layer and wall-shear prediction &
Evaluate $y^+$ and wall-normal grid distribution. \\

\addlinespace

Turbulence closure &
Changed shock position or separation behavior &
Compare suitable models with experimental evidence. \\

\addlinespace

Discretization scheme &
Numerical diffusion or numerical oscillation &
Evaluate scheme sensitivity and grid refinement. \\

\addlinespace

Boundary conditions &
Incorrect operating point or artificial domain influence &
Verify freestream quantities and test domain sensitivity. \\

\addlinespace

Convergence &
Drifting pressure or aerodynamic coefficients &
Monitor residuals and engineering quantities simultaneously. \\

\addlinespace

Physical unsteadiness &
Time-varying shock position and aerodynamic forces &
Use time-accurate simulation and statistical averaging. \\

\bottomrule

\end{tabularx}

\end{table}

% ============================================================
% SYNTHESIS
% ============================================================

\subsection{Synthesis and Remaining Research Questions}

The studies reviewed demonstrate that numerical accuracy cannot be
reduced to a single CFD setting.

Mesh refinement is necessary for assessing discretization sensitivity,
but a grid-independent calculation may still disagree with experiment
because of turbulence-model or boundary-condition errors.

Similarly, agreement between two turbulence models demonstrates
numerical consistency but does not by itself establish physical
accuracy.

Integrated aerodynamic coefficients can also conceal local flow errors.
A CFD model may predict a reasonable lift coefficient while placing the
shock at an incorrect streamwise position.

For transonic validation, the most convincing evidence therefore comes
from several quantities considered together: pressure coefficient,
shock position, aerodynamic forces, grid sensitivity, and numerical
convergence.

Several research challenges remain.

Transition between laminar and turbulent boundary layers can influence
shock behavior but is often simplified in engineering RANS calculations.

Three-dimensional effects also limit the direct application of
two-dimensional airfoil results to real wings. Practical configurations
include finite span, sweep, crossflow, and more complex shock geometry.

Unsteady buffet creates another challenge because a moving physical
shock should not necessarily converge to one constant location. URANS,
LES, and hybrid approaches can provide additional information but
increase computational cost.

For shock-control technology, robust optimization across several Mach
numbers and angles of attack may be more useful than optimization for a
single operating point.

% ============================================================
% CONCLUSION
% ============================================================

\section{Conclusion}

Transonic-airfoil CFD is challenging because localized supersonic
regions, shock waves, turbulent boundary layers, separation, and
possible unsteadiness can occur within the same flow field.

The literature reviewed here shows that numerical reliability depends
on several interconnected choices.

Mesh resolution affects the representation of the boundary layer and
shock, but cell count alone does not determine accuracy. Turbulence
closure can alter shock position and separation behavior, while
numerical convergence does not necessarily imply agreement with
physical measurements.

The RAE 2822 remains valuable because benchmark experimental data allow
researchers to compare numerical predictions with independent evidence.

Comparisons of surface pressure, shock position, aerodynamic
coefficients, grid sensitivity, and convergence behavior provide a
stronger validation basis than any individual metric.

Shock-control bumps demonstrate why this process matters beyond
benchmarking.

Because bump geometry and placement depend on the baseline shock
structure, numerical errors in shock prediction can propagate directly
into aerodynamic optimization.

The main conclusion of this review is therefore that validation should
be treated as part of the predictive engineering process rather than a
final procedural step.

Reliable transonic CFD depends not on how convincing the resulting flow
visualization appears, but on whether multiple numerical and
experimental indicators support the physical accuracy of the prediction.

% ============================================================
% AI ACKNOWLEDGEMENT
% Required by current IJHSR template
% ============================================================

\section{AI Acknowledgement}

During the preparation of this demonstration manuscript, the author
used OpenAI ChatGPT to assist with document formatting, organization,
and creation of a LaTeX representation of the IJHSR Review Article
Template. The resulting content was reviewed and edited by the author.
This document is intended as a personal demonstration of journal
formatting and is not intended for submission.

% ============================================================
% ACKNOWLEDGEMENTS
% ============================================================

\section{Acknowledgements}

The author acknowledges the University of North Texas engineering
environment for providing exposure to computational aerodynamics and
transonic CFD concepts.

% ============================================================
% REFERENCES
% ACS style
% Numbers follow first appearance
% ============================================================

\section{References}

\begin{enumerate}[
    label=\arabic*.,
    leftmargin=1.7em,
    itemsep=4pt,
    topsep=0pt
]

\item
Oberkampf, W. L.; Trucano, T. G.
Verification and Validation in Computational Fluid Dynamics.
\textit{Prog. Aerosp. Sci.}
\textbf{2002}, \textit{38}, 209--272.
\url{https://doi.org/10.1016/S0376-0421(02)00005-2}.

\item
Cook, P. H.; McDonald, M. A.; Firmin, M. C. P.
\textit{Aerofoil RAE 2822: Pressure Distributions, and Boundary Layer
and Wake Measurements};
AGARD Advisory Report AR-138, 1979.

\item
Liou, W. W.; Huang, P. G.; Shih, T.-H.
Turbulence Model Assessment for Shock Wave/Turbulent Boundary-Layer
Interaction in Transonic and Supersonic Flows.
\textit{Comput. Fluids}
\textbf{2000}, \textit{29}, 275--299.
\url{https://doi.org/10.1016/S0045-7930(99)00010-9}.

\item
Celik, I. B.; Ghia, U.; Roache, P. J.; Freitas, C. J.;
Coleman, H.; Raad, P. E.
Procedure for Estimation and Reporting of Uncertainty Due to
Discretization in CFD Applications.
\textit{J. Fluids Eng.}
\textbf{2008}, \textit{130}, 078001.
\url{https://doi.org/10.1115/1.2960953}.

\item
NASA NPARC Alliance Validation Archive.
RAE 2822 Transonic Airfoil: Study No. 5.
NASA Glenn Research Center.

\item
Malkocoglu, U. E.; Barada, M. E.-H. A.; Cetin, F.;
Celik, B.; Yuceil, K. B.
Experimental and Computational Investigations of Shock Wave--Boundary
Layer Interaction on a Transonic Airfoil with and without Bump.
\textit{AIAA AVIATION Forum}, 2020, AIAA Paper 2020-2759.
\url{https://doi.org/10.2514/6.2020-2759}.

\item
Spalart, P. R.; Allmaras, S. R.
A One-Equation Turbulence Model for Aerodynamic Flows.
\textit{30th AIAA Aerospace Sciences Meeting and Exhibit},
1992, AIAA Paper 92-0439.
\url{https://doi.org/10.2514/6.1992-439}.

\item
Menter, F. R.
Two-Equation Eddy-Viscosity Turbulence Models for Engineering
Applications.
\textit{AIAA J.}
\textbf{1994}, \textit{32}, 1598--1605.
\url{https://doi.org/10.2514/3.12149}.

\item
Bruce, P. J. K.; Colliss, S. P.
Review of Research into Shock Control Bumps.
\textit{Shock Waves}
\textbf{2015}, \textit{25}, 451--471.
\url{https://doi.org/10.1007/s00193-014-0533-4}.

\item
Tian, Y.; Gao, S.; Liu, P.; Wang, J.
Transonic Buffet Control Research with Two Types of Shock Control
Bump Based on RAE2822 Airfoil.
\textit{Chin. J. Aeronaut.}
\textbf{2017}, \textit{30}, 1681--1696.
\url{https://doi.org/10.1016/j.cja.2017.07.011}.

\end{enumerate}

% ============================================================
% AUTHORS
% Optional but recommended
% Maximum 50 words
% ============================================================

\section{Authors}

\textbf{Adhvaith Ananth} is a student at the Texas Academy of
Mathematics and Science at the University of North Texas. His academic
interests include aerospace engineering, computational fluid dynamics,
transonic aerodynamics, and high-speed flow analysis.

\end{document}